\section{Objective and risk preference}
\label{sec:objective}

Our objective was to outperform the S\&P 500 Total Return benchmark through
systematic long-short stock selection. The competition ran from
13 February to 6 May 2026 with \$1 million, S\&P 500 stocks, at least ten longs
and ten shorts, and a 200\% gross exposure ceiling. Our simulations apply its
2\% cash and 8\% borrowing annual percentage rates, plus a \$2 fee per order
and the bid-ask spread.

We calibrate risk preference from ten questionnaire responses, reversing
items 2, 3, 8 and 9 as \(Q_i^*=8-Q_i\):
\begin{center}
\small
\begin{tabular}{lrrrrrrrrrr}
\toprule
Item & 1 & 2 & 3 & 4 & 5 & 6 & 7 & 8 & 9 & 10\\
\midrule
Our response & 7 & 5 & 1 & 1 & 7 & 5 & 7 & 1 & 3 & 3\\
Adjusted score & 7 & 3 & 7 & 1 & 7 & 5 & 7 & 7 & 5 & 3\\
\bottomrule
\end{tabular}
\end{center}
\vspace{-0.4em}
\begin{equation}
 \mathrm{BRTI}=\frac{52}{10}=5.2,\qquad
 q=\frac{5.2-1}{6}=0.7,\qquad
 \lambda=12.5e^{-0.35(5.2-1)}=2.8741.
\end{equation}

This indicates relatively high risk tolerance. The frontier analysis uses
\(U=\mathbb{E}[R_p]-(\lambda/2)\sigma_p^2\), with annual decimal returns
and variance; the calibration does not set strategy leverage.
